\chapter{I thread}

Il server sa già tenere più client insieme: una \texttt{fork} per connessione, un figlio ciascuno. Ma quando un client
scrive una riga, come arriva agli altri due? \textbf{Non arriva}. Alla \texttt{fork} ogni figlio ha ricevuto una copia
privata della memoria, quindi non esiste «la lista dei client connessi».
\begin{itemize}
    \item Si potrebbe mandare la riga da un processo all'altro, ma sarebbe un secondo programma, con un suo protocollo, che
          non abbiamo scritto.
    \item Ogni connessione costa inoltre uno spazio di indirizzamento, le sue tabelle delle pagine e una nuova voce nelle
          tabelle dello scheduler.
    \item Quello che vogliamo è un altro \textbf{flusso di controllo dentro lo stesso processo}, che condivida tutto per
          impostazione predefinita.
\end{itemize}

% ══════════════════════════════════════════════════════════════
\section{Un processo, più flussi di controllo}
% ══════════════════════════════════════════════════════════════

\dfn{Thread}{
    Un \textbf{thread} è un flusso di controllo dentro un processo: ha il proprio stack e i propri registri, e tutto il resto
    in comune con gli altri thread dello stesso processo.
}

\begin{figure}[H]
\centering
\begin{tikzpicture}[every node/.style={font=\scriptsize\sffamily},
    sh/.style={draw=netblue!70!black, fill=netblue!8, minimum width=7.8cm, minimum height=0.5cm, align=center},
    th/.style={draw=netgreen!70!black, fill=netgreen!10, rounded corners=2pt, minimum width=2.3cm, align=center}]
  \node[draw=black!60, thick, rounded corners, minimum width=8.6cm, minimum height=5.6cm] (p) at (0,-0.25) {};
  \node[anchor=north west, font=\small\bfseries\sffamily] at (p.north west) {\ processo};
  \node[sh] at (0,1.65) {codice};
  \node[sh] at (0,1.1)  {dati globali};
  \node[sh] at (0,0.55) {heap};
  \node[sh] at (0,0.0)  {tabella dei descriptor: socket e file};
  \foreach \x/\n in {-2.75/1, 0/2, 2.75/3} {
    \node[th] at (\x,-1.1) {registri\\program counter};
    \node[th] at (\x,-1.9) {stack};
    \draw[-{Stealth[length=2mm]}, thick, netred] (\x+1.3,-0.6) -- (\x+1.3,-2.3);
    \node[text=netred] at (\x,-2.55) {thread \n};
  }
\end{tikzpicture}
\caption{Il processo non è più ciò che gira: è il \textbf{contenitore}. Lo scheduler esegue un thread, e fino a oggi ogni
         processo ne aveva uno solo.}
\end{figure}

\subsection{Che cosa è privato e che cosa è condiviso}

\begin{table}[H]
\centering
\begin{tabular}{l l}
\toprule
\textbf{Privato di ogni thread} & \textbf{Condiviso da tutti i thread del processo} \\
\midrule
thread id, registri, program counter & il codice, e ogni variabile globale \\
il suo stack, quindi le sue variabili locali & lo heap: tutto ciò che ha restituito \texttt{malloc} \\
la maschera dei segnali & la tabella dei descriptor: socket e file \\
\texttt{errno} & la disposizione di ogni segnale (gli handler) \\
la priorità di scheduling & pid, ppid, directory di lavoro, credenziali \\
\bottomrule
\end{tabular}
\caption{La colonna di sinistra è corta di proposito: un thread possiede lo stato di \textbf{una computazione}, nient'altro.}
\end{table}

In particolare, un descriptor aperto da un thread è \texttt{fd 4} \textbf{per tutti}.

\begin{figure}[H]
\centering
\begin{tikzpicture}[every node/.style={font=\scriptsize\sffamily},
    th/.style={draw=netgreen!70!black, fill=netgreen!10, rounded corners=2pt, minimum width=2.4cm, align=center}]
  \node[draw=black!60, thick, rounded corners, minimum width=9.2cm, minimum height=4.0cm] (p) at (0,0) {};
  \node[anchor=north, font=\small\bfseries\sffamily] at (p.north) {un processo, uno spazio di indirizzamento};
  \foreach \x/\n/\f/\c in {-3/1/4/A, 0/2/5/B, 3/3/6/C}
    \node[th] at (\x,0.75) {\textbf{thread \n}\\stack proprio\\\texttt{fd \f} client \c};
  \node[draw=netblue!70!black, fill=netblue!8, rounded corners=2pt, minimum width=8.4cm, align=center] at (0,-0.55)
    {\textbf{condiviso}: lo stato della chat, chi è connesso, che cosa inviargli\\una sola copia, raggiunta leggendo una variabile};
  \node[draw=netred, fill=netred!8, rounded corners=2pt, minimum width=8.4cm, align=center] at (0,-1.45)
    {\textbf{thread principale}: \texttt{accept} sul \texttt{fd 3}, un thread nuovo per ogni connessione accettata};
\end{tikzpicture}
\caption{Una sola tabella dei descriptor, quindi una sola \texttt{close} chiude una connessione; e lo stato della chat è
         una \textbf{variabile}, non un messaggio che qualcuno deve inviare.}
\end{figure}

\subsection{Che cosa compra la condivisione, e che cosa costa}

\begin{itemize}
    \item \textbf{Niente comunicazione fra processi}: la lista dei client connessi è una variabile; un thread la scrive, gli
          altri la leggono, e niente viene trasmesso.
    \item \textbf{Più economico}: creare un thread non copia alcuno spazio di indirizzamento, e passare da un thread all'altro
          non cambia le tabelle delle pagine.
    \item \textbf{Blocco indipendente}: un thread fermo in \texttt{read} blocca solo se stesso, purché il kernel sappia che
          esiste.
    \item \textbf{Il prezzo}: niente li separa. Un thread che corrompe una struttura condivisa la corrompe per tutti.
    \item E un thread che dereferenzia un puntatore nullo porta giù l'intero server, client compresi; un processo figlio
          portava giù soltanto se stesso.
\end{itemize}

Un server con \texttt{fork} è isolato per costruzione; uno con i thread no, ed è di questo che parla la seconda metà del
capitolo.

\subsection{Concorrenza non è parallelismo}

\begin{figure}[H]
\centering
\begin{tikzpicture}[every node/.style={font=\scriptsize\sffamily},
    t/.style={draw=black!50, minimum height=0.4cm, anchor=west, inner sep=0pt, fill=#1}]
  \node[anchor=west, font=\small\bfseries\sffamily] at (-2.6,0.8) {Concorrenza: un core};
  \node[anchor=east] at (-0.2,0) {core};
  \foreach \x/\c in {0/netblue!40, 1.0/netgreen!40, 2.0/netorange!40, 3.0/netblue!40, 4.0/netgreen!40, 5.0/netorange!40}
    \node[t=\c, minimum width=1.0cm] at (\x,0) {};
  \node[anchor=west, font=\small\bfseries\sffamily] at (-2.6,-0.8) {Parallelismo: tre core};
  \foreach \y/\c/\n in {-1.5/netblue!40/1, -2.0/netgreen!40/2, -2.5/netorange!40/3} {
    \node[anchor=east] at (-0.2,\y) {core \n};
    \node[t=\c, minimum width=6.0cm] at (0,\y) {};
  }
  \draw[-{Stealth[length=2mm]}, black!60] (0,-3.0) -- (6.4,-3.0) node[right] {tempo};
\end{tikzpicture}
\caption{\textbf{Concorrenza}: più compiti in corso insieme, alternati. \textbf{Parallelismo}: più compiti eseguiti nello
         stesso istante.}
\end{figure}

A un server serve la prima: i suoi thread passano la vita ad \textbf{aspettare la rete}, non a calcolare. La seconda è un
regalo che l'hardware può fare.

% ══════════════════════════════════════════════════════════════
\section{Dove sono conosciuti i thread}
% ══════════════════════════════════════════════════════════════

\begin{figure}[H]
\centering
\begin{tikzpicture}[every node/.style={font=\scriptsize\sffamily},
    ut/.style={draw=netgreen!70!black, fill=netgreen!12, circle, minimum size=0.5cm, inner sep=0pt},
    kt/.style={draw=netblue!70!black, fill=netblue!15, circle, minimum size=0.5cm, inner sep=0pt}]
  % livello utente
  \begin{scope}
    \node[font=\small\bfseries\sffamily] at (1.5,2.3) {Thread a livello utente};
    \draw[black!50, dashed] (-0.6,0.6) -- (3.6,0.6);
    \node[anchor=west, text=black!60] at (3.6,1.2) {spazio utente};
    \node[anchor=west, text=black!60] at (3.6,0.0) {kernel};
    \foreach \x in {0.5,1.5,2.5} \node[ut] at (\x,1.5) {};
    \node[draw, rounded corners=2pt, fill=lapp!40] (lib) at (1.5,0.95) {libreria dei thread};
    \node[kt] (k) at (1.5,0.0) {};
    \draw (lib) -- (k);
    \node[text=netred, align=center] at (1.5,-0.7) {il kernel vede\\una sola cosa eseguibile};
  \end{scope}
  % livello kernel
  \begin{scope}[xshift=6.6cm]
    \node[font=\small\bfseries\sffamily] at (1.5,2.3) {Thread a livello kernel};
    \draw[black!50, dashed] (-0.6,0.6) -- (3.6,0.6);
    \foreach \x in {0.5,1.5,2.5} { \node[ut] at (\x,1.5) {}; \node[kt] at (\x,0.0) {}; \draw (\x,1.25) -- (\x,0.25); }
    \node[draw, rounded corners=2pt, fill=lnet!40] at (1.5,-0.55) {tabella dei thread del kernel};
    \node[text=netgreen!60!black, align=center] at (1.5,-1.15) {ognuno schedulato\\e bloccato a sé};
  \end{scope}
\end{tikzpicture}
\caption{A sinistra la libreria fa tutto nello spazio utente: cambiare thread è una chiamata di funzione, ma il kernel vede
         una sola cosa eseguibile. A destra il kernel ha una propria tabella: schedula ogni thread e blocca ciascuno
         separatamente.}
\end{figure}

\begin{itemize}
    \item \textbf{Molti a uno}: tutti i thread utente su un solo thread del kernel. Cambiare è economico, ma una chiamata
          bloccante li ferma tutti, e due core non vengono mai usati.
    \item \textbf{Uno a uno}: un thread del kernel per ogni thread utente: blocco vero, parallelismo vero, e un oggetto del
          kernel per ciascuno. È \textbf{Linux}, e quindi il nostro caso.
    \item \textbf{Molti a molti}: $m$ thread utente schedulati su $n$ thread del kernel. È la risposta generale, abbastanza
          complicata da far tornare indietro i sistemi che l'avevano provata.
\end{itemize}
Nel nostro server ogni thread viene quindi bloccato e schedulato per conto suo, e ogni thread è un oggetto del kernel: per
questo non se ne possono avere un milione.

% ══════════════════════════════════════════════════════════════
\section{Creare un thread}
% ══════════════════════════════════════════════════════════════

\begin{lstlisting}[language={[POSIX]C}]
#include <pthread.h>

int pthread_create(pthread_t *tid, const pthread_attr_t *attr,
                   void *(*start_routine)(void *), void *arg);
\end{lstlisting}

\begin{itemize}
    \item \texttt{tid}: un \textbf{output}, l'identità del nuovo thread, di tipo opaco \texttt{pthread\_t}. Un thread chiede
          la propria con \texttt{pthread\_self()}.
    \item \texttt{attr}: gli attributi, o \texttt{NULL} per quelli predefiniti, che vanno bene finché non serve un thread
          \emph{detached}.
    \item \texttt{start\_routine}: dove comincia il nuovo thread. Non prosegue da qui: \textbf{chiama} quella funzione.
    \item \texttt{arg}: l'unico argomento che quella funzione riceve.
\end{itemize}
Si compila con \texttt{gcc -pthread}: collega la libreria e imposta i flag che gli header si aspettano.

\warning{Gli errori non passano per errno}{
    Tutte le chiamate viste finora restituivano $-1$ e impostavano \texttt{errno}. Le funzioni \texttt{pthread}
    \textbf{restituiscono il codice d'errore} e lasciano stare \texttt{errno}: si stampa con \texttt{strerror(err)}, non
    con \texttt{perror}, che stamperebbe l'ultimo fallimento non correlato. Il motivo è storico (una globale sarebbe
    ambigua con i thread) e oggi \texttt{errno} è in effetti per thread. \texttt{pthread\_create} fallisce davvero:
    \texttt{EAGAIN} quando il sistema non concede un altro thread, e un server deve sopravvivere.
}

\begin{lstlisting}[language={[POSIX]C}, style=wnumbers]
int err = pthread_create(&tid, NULL, worker, arg);
if (err != 0)
    fprintf(stderr, "pthread_create: %s\n", strerror(err));   /* non perror */
\end{lstlisting}

\subsection{La funzione di partenza e il suo unico argomento}

\begin{lstlisting}[language={[POSIX]C}, style=wnumbers, title={Un argomento solo: una struct per passarne di piu'}]
struct job { int fd; unsigned long n; };     /* cio' che serve a un thread */

static void *worker(void *arg) {
    struct job *j = arg;                     /* il void * torna com'era partito */
    serve(j->fd, j->n);
    close(j->fd);
    free(j);                                 /* e' di questo thread, quindi lo libera lui */
    return NULL;                             /* come pthread_exit(NULL) */
}

struct job *j = malloc(sizeof *j);
j->fd = fd; j->n = ++served;
pthread_create(&tid, NULL, worker, j);       /* un puntatore, e qualunque cosa dietro */
\end{lstlisting}

Più di un valore vuol dire una struttura, e quella struttura deve \textbf{esistere ancora} quando il thread la legge.

\begin{warningbox}{{L'argomento non dev'essere una variabile riusata}}
\begin{lstlisting}[language={[POSIX]C}]
for (;;) {
    int fd = accept(listen_fd, NULL, NULL);
    pthread_create(&tid, NULL, worker, &fd);    /* una sola variabile, per tutti */
}
\end{lstlisting}
Due client si connettono a un millisecondo di distanza: quale descriptor serve ciascun thread? Quello che c'è in
\texttt{fd} \textbf{quando guarda}. La seconda \texttt{accept} può averlo già sovrascritto: due thread sulla stessa
connessione, e una connessione che nessuno serve. Si dà a ogni thread \textbf{un oggetto suo}: una \texttt{malloc} per
connessione, liberata dal thread che la riceve. Nei test funziona, perché nei test i client arrivano uno alla volta.
\end{warningbox}

\subsection{Due thread e un main}

\begin{lstlisting}[language={[POSIX]C}, style=wnumbers, title={hello.c}]
static void *say(void *arg) {
    printf("%s: tid %lu\n", (char *) arg, (unsigned long) pthread_self());
    return NULL;
}

int main(void) {
    pthread_t t1, t2;
    pthread_create(&t1, NULL, say, "first");
    pthread_create(&t2, NULL, say, "second");
    pthread_join(t1, NULL);
    pthread_join(t2, NULL);
    return 0;
}
\end{lstlisting}

Nulla dice quale dei due stampi per primo, ed eseguirlo due volte può dare risposte diverse. \texttt{pthread\_t} è
opaco: stamparlo è una comodità di questo sistema, non una cosa portabile.

\warning{Quando main ritorna, muoiono tutti}{
    Togliendo le due \texttt{join}, che cosa stampa il programma? Spesso \textbf{niente}. Il \texttt{return} da
    \texttt{main} è una \texttt{exit}, e \texttt{exit} termina il processo, compresi i thread che stavano per stampare.
    \begin{itemize}
        \item \texttt{exit} chiamata da \textbf{qualunque} thread termina tutti i thread: non esiste una sua versione che
              termini solo il chiamante.
        \item O si aspettano con \texttt{join} i thread avviati, oppure, se \texttt{main} non ha altro da fare, lo si chiude con
              \texttt{pthread\_exit(NULL)}: il processo vive allora fino alla fine del suo ultimo thread.
        \item Una \texttt{sleep} in \texttt{main} non è un modo di aspettare un thread: è una scommessa su quanto dura il
              lavoro.
    \end{itemize}
}

% ══════════════════════════════════════════════════════════════
\section{La fine di un thread}
% ══════════════════════════════════════════════════════════════

\begin{itemize}
    \item \textbf{return} dalla funzione di partenza: il valore restituito è il valore di uscita del thread, ed è il modo
          normale.
    \item \textbf{\texttt{pthread\_exit(valore)}}: la stessa cosa da qualunque punto, anche da una funzione chiamata dieci
          livelli più in basso.
    \item \textbf{Cancellato}: un altro thread ne ha chiesto la terminazione; il valore di uscita è allora
          \texttt{PTHREAD\_CANCELED}.
    \item \textbf{Non sono modi di finire un thread}: \texttt{exit}, \texttt{\_exit}, o un segnale la cui azione è
          terminare: quelli chiudono l'intero processo.
\end{itemize}
In qualunque modo finisca, il valore lasciato deve \textbf{sopravvivere al thread}: il suo stack non lo fa.

\subsection{pthread\_join}

\begin{lstlisting}[language={[POSIX]C}]
int pthread_join(pthread_t tid, void **retval);

void *status;
pthread_join(tid, &status);      /* blocca finche' quel thread termina */
pthread_join(tid, NULL);         /* il valore di uscita non ci interessa */
\end{lstlisting}

\begin{itemize}
    \item Aspetta \textbf{un thread con un nome}: non esiste \texttt{pthread\_join(-1)}, un «chiunque finisca per primo»
          come con \texttt{wait}.
    \item Se quel thread è già terminato ritorna subito, esattamente come aspettare un figlio già zombie.
    \item Ed è l'analogia da tenere: un thread terminato e non aspettato occupa ancora le sue risorse.
\end{itemize}

\begin{lstlisting}[language={[POSIX]C}, style=wnumbers, title={Il valore deve sopravvivere al thread}]
static void *local(void *arg) {
    struct pair p = { 1, 2 };            /* sullo stack di questo thread */
    return &p;                           /* che fra un attimo non c'e' piu' */
}
static void *heap(void *arg) {
    struct pair *p = malloc(sizeof *p);
    p->a = 5; p->b = 6;
    return p;                            /* chi fa la join lo legge, poi lo libera */
}
\end{lstlisting}

\begin{lstlisting}[language=bash]
thread 1: a=1 b=2
master  : a=1075156600 b=1077281896      <- lo stack nel frattempo riusato
thread 3: a=5 b=6
master  : a=5 b=6
\end{lstlisting}

Una globale o lo heap; mai lo stack. Il compilatore non lo impedisce, e la prima esecuzione spesso stampa i numeri giusti.

\subsection{Thread detached}

Un thread che nessuno aspetterà mai dovrebbe dirlo, e allora rilascia tutto nel momento in cui termina.

\begin{lstlisting}[language={[POSIX]C}, style=wnumbers]
pthread_detach(tid);                     /* da un altro thread, oppure... */
pthread_detach(pthread_self());          /* ...dal thread stesso, per prima cosa */

pthread_attr_t attr;                     /* oppure deciso alla creazione */
pthread_attr_init(&attr);
pthread_attr_setdetachstate(&attr, PTHREAD_CREATE_DETACHED);
pthread_create(&tid, &attr, worker, j);
pthread_attr_destroy(&attr);
\end{lstlisting}

\begin{itemize}
    \item Per default i thread sono \emph{joinable}: un server che se ne dimentica perde, per ogni client servito, la
          contabilità di un thread.
    \item \emph{Detached} è in pratica irreversibile: da lì in poi \texttt{pthread\_join} su quel thread è un errore, e il suo
          valore di uscita è perso.
\end{itemize}

% ══════════════════════════════════════════════════════════════
\section{Un thread per connessione}
% ══════════════════════════════════════════════════════════════

\begin{lstlisting}[language={[POSIX]C}, style=wnumbers, title={Il ciclo di accept, in sintesi}]
for (;;) {
    int *fd = malloc(sizeof *fd);          /* uno per connessione, liberato dal thread */
    *fd = accept(listen_fd, NULL, NULL);
    if (*fd < 0) { perror("accept"); free(fd); continue; }

    pthread_t tid;
    int err = pthread_create(&tid, NULL, worker, fd);
    if (err != 0) {                        /* finiti i thread: si rifiuta questo client */
        fprintf(stderr, "pthread_create: %s\n", strerror(err));
        close(*fd); free(fd); continue;
    }
    pthread_detach(tid);                   /* nessuno la aspettera' mai */
}
\end{lstlisting}

Si noti che cosa manca: \textbf{nessuna \texttt{close}} nel thread che accetta. Il descriptor è uno solo, non due, quindi
chi lo chiude chiude la connessione.

\begin{lstlisting}[language=bash, title={I thread visti da fuori}]
$ gcc -pthread -o chatd chatd.c
$ ./chatd 5000 &
$ ps -o pid,lwp,nlwp,stat,comm -L -p $(pgrep chatd)
  PID   LWP NLWP STAT COMMAND
 5310  5310    4 Sl   chatd       <- il thread principale, in accept
 5310  5311    4 Sl   chatd
 5310  5312    4 Sl   chatd       <- uno per client connesso
 5310  5313    4 Sl   chatd
\end{lstlisting}

Un pid, quattro LWP (\emph{lightweight process}): il kernel li schedula uno per uno, e \texttt{top -H} li mostra consumare
CPU separatamente. \texttt{kill} invece si rivolge ancora al processo, perché un segnale si manda a un pid.

\begin{table}[H]
\centering
\begin{tabular}{>{\ttfamily}l >{\ttfamily}l l}
\toprule
\normalfont\textbf{Processi} & \normalfont\textbf{Thread} & \textbf{E la differenza} \\
\midrule
fork    & pthread\_create & il figlio riparte da qui, il thread parte in una funzione \\
exit    & pthread\_exit   & \texttt{exit} in un thread termina il processo \\
waitpid & pthread\_join   & niente «uno qualsiasi»: si nomina il thread \\
kill    & pthread\_kill   & la disposizione è una per processo \\
getpid  & pthread\_self   & l'id è opaco, solo confrontabile \\
\normalfont zombie & \normalfont un thread mai aspettato & stessa perdita, stessa cura \\
\bottomrule
\end{tabular}
\caption{Processi e thread, fianco a fianco. I nomi si corrispondono, ed è una buona mappa, salvo che i processi sono isolati
         e i thread no.}
\end{table}

% ══════════════════════════════════════════════════════════════
\section{Il prezzo della condivisione}
% ══════════════════════════════════════════════════════════════

\subsection{Il contatore che perde gli incrementi}

\begin{lstlisting}[language={[POSIX]C}, style=wnumbers, title={bump.c}]
static int count = 0;                    /* una variabile, tutti i thread */

static void *bump(void *arg) {
    for (int i = 0; i < 100000; i++)
        count++;
    return NULL;
}
/* main: ne crea otto, li aspetta tutti e otto, stampa count */
\end{lstlisting}

\begin{lstlisting}[language=bash]
$ ./bump ; ./bump ; ./bump
count = 431027      <- atteso 800000
count = 800000      <- a volte giusto: ancora peggio
count = 512663
\end{lstlisting}

Nessun thread sbaglia e nessuna riga di C è sbagliata. E non è una questione di due core: su un core solo gli stessi
incrementi si perdono, perché un thread può essere interrotto \textbf{nel mezzo} di un incremento.

\begin{table}[H]
\centering
\begin{tabular}{l l c}
\toprule
\textbf{Thread A} & \textbf{Thread B} & \texttt{count} \\
\midrule
load $\rightarrow$ 7 &                      & 7 \\
                     & load $\rightarrow$ 7 & 7 \\
add, store 8         &                      & 8 \\
                     & add, store 8         & 8 \\
\bottomrule
\end{tabular}
\caption{\texttt{count++} è un carattere solo ma tre operazioni macchina: load, add, store. Due incrementi, uno perso.}
\end{table}

Ogni volta che due thread possono toccare la stessa variabile e almeno uno la scrive, il risultato dipende dai tempi: è
una \key{race condition}, e serve un modo per vietare quella sovrapposizione. È l'argomento del prossimo capitolo.

\subsection{Non tutte le funzioni possono essere chiamate da due thread}

Le nostre globali non sono l'unico stato condiviso: alcune funzioni di libreria tengono il proprio, in un buffer statico
che non si vede.
\begin{itemize}
    \item \texttt{strtok} ricorda dove si era fermata: due thread che analizzano due righe si calpestano la posizione.
    \item \texttt{gethostbyname}, \texttt{ctime}, \texttt{inet\_ntoa} restituiscono un puntatore a \textbf{un solo} buffer,
          sovrascritto dal chiamante successivo.
    \item Alcune hanno un gemello \textbf{rientrante} che prende il buffer da noi: \texttt{strtok\_r}, \texttt{rand\_r}. Si
          preferiscono loro, e si preferisce \texttt{getaddrinfo}, progettata così fin dall'inizio.
    \item \texttt{errno} non è fra queste: è per thread, quindi il fallimento di un altro non compare mai nel nostro.
\end{itemize}
La pagina del manuale dice \emph{MT-Safe} oppure no, e vale la pena guardarla prima di mettere una funzione dentro un worker.

\subsection{Thread e segnali}

\begin{itemize}
    \item La \textbf{disposizione} è una per processo: una \texttt{sigaction}, e vale per tutti i thread. Non esiste un
          handler per thread.
    \item La \textbf{maschera} è una per thread: ciascuno decide che cosa bloccare, con \texttt{pthread\_sigmask}.
    \item Un segnale inviato al processo viene consegnato a \textbf{un} thread qualsiasi che non lo blocchi, e non si può
          sapere quale.
    \item \texttt{pthread\_kill(tid, sig)} nomina un thread; ma se la disposizione è terminare, muore comunque tutto il
          processo.
\end{itemize}
Per questo un server con i thread di solito blocca i segnali ovunque e lascia che sia un thread solo ad aspettarli; e
ignora \texttt{SIGPIPE} con \texttt{signal(SIGPIPE, SIG\_IGN)}, così la scrittura su una connessione chiusa restituisce
\texttt{EPIPE} invece di uccidere il server.

\warning{Un thread chiama fork}{
    Uno di otto thread chiama \texttt{fork}: quanti thread ha il figlio? \textbf{Uno}, quello che l'ha chiamata. Gli altri
    sette non esistono lì, ma tutto quello che avevano fatto alla memoria sì.
    \begin{itemize}
        \item Quello che avevano in mano resta com'era: una struttura aggiornata a metà resta a metà, e un lock resta chiuso
              per sempre, perché il suo proprietario non c'è a rilasciarlo.
        \item Il figlio di un processo con più thread può quindi fare solo le cose sicure: in pratica, una \texttt{exec},
              subito.
        \item \texttt{exec} non ne risente: sostituisce il programma, e il nuovo parte con il solo thread che l'ha chiamata.
    \end{itemize}
}

\subsection{Cancellare un thread, e perché un server lo fa di rado}

\begin{lstlisting}[language={[POSIX]C}]
pthread_cancel(tid);              /* chiede la terminazione; non aspetta */
pthread_join(tid, &status);       /* status == PTHREAD_CANCELED */
\end{lstlisting}

\begin{itemize}
    \item Per default la cancellazione è \emph{differita}: ha effetto a un \textbf{punto di cancellazione} (\texttt{read},
          \texttt{write}, \texttt{accept}, \texttt{sleep}), quindi un thread occupato a calcolare non se ne accorge per un
          po'.
    \item Il thread si ferma dove si trova: il suo buffer, il suo descriptor e qualunque cosa avesse bloccato restano com'erano.
    \item Per questo un server \textbf{chiede} invece di ordinare: un flag \texttt{volatile sig\_atomic\_t} controllato dai
          worker fra una richiesta e l'altra, oppure semplicemente chiude il descriptor e lascia che la loro \texttt{read}
          ritorni.
\end{itemize}

\subsection{Quanto costa un thread, e quanti se ne possono avere}

\begin{itemize}
    \item Crearne uno costa all'incirca \textbf{un ordine di grandezza meno} di una \texttt{fork}: nessuno spazio di
          indirizzamento da duplicare, nessuna tabella delle pagine da copiare.
    \item Passare fra due thread dello stesso processo non cambia lo spazio di indirizzamento, quindi le cache e il TLB
          sopravvivono.
    \item Ma ogni thread riserva uno stack: su Linux 8 MB di spazio di indirizzamento per default, e se ne può chiedere
          meno alla creazione.
    \item E ognuno è un oggetto del kernel con il suo stack del kernel: migliaia sono normali, centinaia di migliaia no.
\end{itemize}
Un thread per client va bene per una chat con un'aula collegata, non per diecimila connessioni: lì subentrano un
\textbf{insieme fisso di worker} e un thread che osserva \textbf{molte connessioni insieme}.

% ══════════════════════════════════════════════════════════════
\section{Esercizi}
% ══════════════════════════════════════════════════════════════

\begin{esercizio}{Contatori e thread}
Scrivere due programmi brevi.
\begin{enumerate}[label=(\alph*)]
    \item Prendere $n$ dalla riga di comando, creare $n$ thread che dormono ciascuno un numero casuale di secondi, poi
          incrementano un contatore globale e lo stampano; aspettarli tutti.
    \item Lo stesso, senza \texttt{sleep}, con ogni thread che incrementa 100\,000 volte.
\end{enumerate}
Eseguire ciascuno venti volte e contare le risposte. Ci sono race condition nel primo programma? E perché il secondo le
mostra mentre il primo le nasconde?

\svolgimento
\begin{lstlisting}[language={[POSIX]C}, style=wnumbers, title={threads.c: la versione (a)}]
static int counter = 0;

static void *work(void *arg) {
    unsigned seed = (unsigned) (uintptr_t) arg;   /* un seme diverso per thread */
    sleep(rand_r(&seed) % 3 + 1);
    counter++;                                    /* la stessa race del capitolo */
    printf("counter = %d\n", counter);
    return NULL;
}

int main(int argc, char *argv[]) {
    int n = atoi(argv[1]);
    pthread_t *t = malloc(n * sizeof *t);
    for (int i = 0; i < n; i++)
        pthread_create(&t[i], NULL, work, (void *) (uintptr_t) (i + 1));
    for (int i = 0; i < n; i++)
        pthread_join(t[i], NULL);
    free(t);
    return 0;
}
\end{lstlisting}
La versione (b) toglie la \texttt{sleep} e mette \texttt{counter++} in un ciclo di 100\,000 iterazioni, stampando
solo alla fine in \texttt{main}. Si usa \texttt{rand\_r} e non \texttt{rand}, che tiene il proprio stato in una globale.

Sì, la race c'è \textbf{anche} nel primo: \texttt{counter++} non è indivisibile, e due thread che si svegliano insieme
possono perdere un incremento. Ma i thread si svegliano a secondi di distanza, e ognuno esegue un solo incremento, che
dura nanosecondi: la probabilità che due di essi si sovrappongano proprio in quelle tre istruzioni è minima, quindi il
risultato è quasi sempre giusto. Nel secondo otto thread eseguono centinaia di migliaia di incrementi negli stessi
millisecondi, e le sovrapposizioni diventano inevitabili. La race non dipende da quanto spesso si vede: un programma
che «funziona sempre» nei test può averla.
\end{esercizio}

\Needspace{20\baselineskip}
% ══════════════════════════════════════════════════════════════
\section{Riepilogo}
% ══════════════════════════════════════════════════════════════

\riepilogo{
\begin{itemize}
    \item Un thread possiede il proprio stack, i propri registri e la propria maschera dei segnali; tutto il resto (globali,
          heap, descriptor) appartiene al processo.
    \item Quindi i thread condividono lo stato \textbf{leggendo una variabile}, che è proprio ciò che un server di chat con i
          processi non poteva fare.
    \item \texttt{pthread\_create} prende una funzione e un argomento, e quell'argomento deve appartenere al thread a cui viene
          dato.
    \item Le chiamate \texttt{pthread} restituiscono un codice d'errore, non $-1$; e un thread terminato va aspettato con
          \texttt{join} o reso \emph{detached}, altrimenti perde risorse.
    \item \texttt{exit} e il ritorno di \texttt{main} terminano tutti i thread; solo \texttt{pthread\_exit} e
          \texttt{return} ne terminano uno.
    \item E il prezzo della condivisione è \texttt{count++}: codice corretto, risultato sbagliato, ogni volta che due thread
          scrivono la stessa cosa.
\end{itemize}
}
